\begin{center}
  \section*{Homework 05 - 3.2, 3.3}
  Due Mar 13 \\
  Uzair Hamed Mohammed
\end{center}

% 10, 12, 14, 16, 20, 22, 27, 28, 30, 36
% 2, 8, 14

\subsection*{3.2 Binomial and Hypergeometric Probabilities}
\begin{enumerate}
  \item[10.] The gunner on a small assault boat fires six missiles at
    an attacking plane. Each has a 20\% chance of being on-target. If
    two or more of the shells find their mark, the plane will crash.
    At the same time, the pilot of the plane fires ten air-to-surface
    rockets, each of which has a 0.05 chance of critically disabling
    the boat. Would you rather be on the plane or the boat?

    \underline{Sol}:\\
    $$
    \begin{aligned}
      & \text{Plane Survival Probability: } \\
      & P(X \ge 2) = 1 - \left[ \binom{6}{0} (0.2)^0 (0.8)^6 +
      \binom{6}{1} (0.2)^1 (0.8)^5 \right] \\
      & P(X \ge 2) = 1 - [0.262144 + 0.393216] = 0.34464 \\
      & P(\text{Survival}_{\text{Plane}}) = 1 - 0.34464 = 0.65536 \\
      \\
      & \text{Boat Survival Probability: } \\
      & P(Y \ge 1) = 1 - \binom{10}{0} (0.05)^0 (0.95)^{10} \\
      & P(Y \ge 1) = 1 - 0.598737 = 0.401263 \\
      & P(\text{Survival}_{\text{Boat}}) = 1 - 0.401263 = 0.598737 \\
      \\
      & P(\text{Survival}_{\text{Plane}}) > P(\text{Survival}_{\text{Boat}}) \\
      & 0.65536 > 0.598737 \\
      \\
      & \boxed{\text{Preference: Plane}}
    \end{aligned}$$

  \item[12.] Experience has shown that only \(\frac{1}{3}\) of all
    patients having a certain disease will recover if given the
    standard treatment. A new drug is to be tested on a group of
    twelve volunteers. If the FDA requires that at least seven of
    these patients recover before it will license the new drug, what
    is the probability that the treatment will be discredited even if
    it has the potential to increase an individual’s recovery rate to
    \(\frac{1}{2}\)?

    \underline{Sol}:\\
    $$
    \begin{aligned}
      & \text{Parameters: } \\
      & n = 12 \\
      & p = 0.5 \\
      \\
      & \text{The drug is discredited if fewer than 7 patients recover: } \\
      & P(X < 7) = P(X \le 6) = \sum_{k=0}^{6} \binom{12}{k} (0.5)^k
      (0.5)^{12-k} \\
      & P(X \le 6) = \frac{1}{2^{12}} \sum_{k=0}^{6} \binom{12}{k} \\
      & \sum_{k=0}^{6} \binom{12}{k} = \binom{12}{0} + \binom{12}{1}
      + \binom{12}{2} + \binom{12}{3} + \binom{12}{4} + \binom{12}{5}
      + \binom{12}{6} \\
      & \sum_{k=0}^{6} \binom{12}{k} = 1 + 12 + 66 + 220 + 495 + 792
      + 924 = 2510 \\
      \\
      & P(X \le 6) = \frac{2510}{4096} \approx \boxed{0.6128}
    \end{aligned}$$

  \item[14.] The captain of a Navy gunboat orders a volley of
    twenty-five missiles to be fired at random along a
    five-hundred-foot stretch of shoreline that he hopes to establish
    as a beachhead. Dug into the beach is a thirty-foot-long bunker
    serving as the enemy’s first line of defense. The captain has
    reason to believe that the bunker will be destroyed if at least
    three of the missiles are on-target. What is the probability of
    that happening?

    \underline{Sol}:\\
    $$
    \begin{aligned}
      & \text{Parameters: } \\
      & n = 25 \\
      & p = \frac{30}{500} = 0.06 \\
      \\
      & \text{Bunker is destroyed if } X \ge 3: \\
      & P(X \ge 3) = 1 - \sum_{k=0}^{2} \binom{25}{k} (0.06)^k (0.94)^{25-k} \\
      & P(X=0) = \binom{25}{0} (0.06)^0 (0.94)^{25} \approx 0.2129 \\
      & P(X=1) = \binom{25}{1} (0.06)^1 (0.94)^{24} \approx 0.3398 \\
      & P(X=2) = \binom{25}{2} (0.06)^2 (0.94)^{23} \approx 0.2602 \\
      \\
      & P(X \ge 3) = 1 - (0.2129 + 0.3398 + 0.2602) \\
      & P(X \ge 3) = 1 - 0.8129 \approx \boxed{0.1871}
    \end{aligned}$$

  \item[16.] Listed in the following table is the length distribution
    of World Series competition for the sixty-four series from 1950
    to 2014 (there was no series in 1994). Assuming that each World
    Series game is an independent event and that the probability of
    either team’s winning any particular contest is 0.5, find the
    probability of each series length. How well does the model fit
    the data? (Compute the “expected” frequencies, that is, multiply
    the probability of a given-length series times 64).

    \underline{Sol}:\\
    $$
    \begin{aligned}
      & \text{Theoretical Probabilities } P(k) \text{ for } k \text{ games:} \\
      & P(k) = 2 \times \binom{k-1}{3} (0.5)^k \\
      & P(4) = 2 \times \binom{3}{3} (0.5)^4 = 2(1)(0.0625) = 0.125 \\
      & P(5) = 2 \times \binom{4}{3} (0.5)^5 = 2(4)(0.03125) = 0.25 \\
      & P(6) = 2 \times \binom{5}{3} (0.5)^6 = 2(10)(0.015625) = 0.3125 \\
      & P(7) = 2 \times \binom{6}{3} (0.5)^7 = 2(20)(0.0078125) = 0.3125 \\
      \\
      & \text{Expected Frequencies } (E_k = P(k) \times 64): \\
      & E_4 = 0.125 \times 64 = 8.0 \\
      & E_5 = 0.25 \times 64 = 16.0 \\
      & E_6 = 0.3125 \times 64 = 20.0 \\
      & E_7 = 0.3125 \times 64 = 20.0 \\
      \\
      & \text{Comparison with Observed Data (1950–2014):} \\
      &
      \begin{array}{c|c|c|c}
        k & \text{Observed } (O_k) & \text{Expected } (E_k) & |O_k - E_k| \\
        \hline
        4 & 12 & 8.0 & 4.0 \\
        5 & 11 & 16.0 & 5.0 \\
        6 & 14 & 20.0 & 6.0 \\
        7 & 27 & 20.0 & 7.0 \\
      \end{array} \\
      \\
      & \text{Model Fit: The model fits poorly. It significantly
      underestimates} \\
      & \text{the frequency of four-game and seven-game series.}
    \end{aligned}$$

  \item[20.] A corporate board contains twelve members. The board
    decides to create a five-person Committee to Hide Corporation
    Debt. Suppose four members of the board are accountants. What is
    the probability that the Committee will contain two accountants
    and three non-accountants?

    \underline{Sol}:\\
    $$
    \begin{aligned}
      & \text{Total members: } N = 12 \\
      & \text{Accountants: } k = 4 \\
      & \text{Non-accountants: } N - k = 8 \\
      & \text{Committee size: } n = 5 \\
      \\
      & P(X = 2) = \frac{\binom{4}{2} \binom{8}{3}}{\binom{12}{5}} \\
      & \binom{4}{2} = \frac{4 \times 3}{2 \times 1} = 6 \\
      & \binom{8}{3} = \frac{8 \times 7 \times 6}{3 \times 2 \times 1} = 56 \\
      & \binom{12}{5} = \frac{12 \times 11 \times 10 \times 9 \times
      8}{5 \times 4 \times 3 \times 2 \times 1} = 792 \\
      \\
      & P(X = 2) = \frac{6 \times 56}{792} = \frac{336}{792} \approx
      \boxed{0.4242}
    \end{aligned}$$

  \item[22.] A city has 4050 children under the age of ten, including
    514 who have not been vaccinated for measles. Sixty-five of the
    city’s children are enrolled in the ABC Day Care Center. Suppose
    the municipal health department sends a doctor and a nurse to ABC
    to immunize any child who has not already been vaccinated. Find a
    formula for the probability that exactly k of the children at ABC
    have not been vaccinated.

    \underline{Sol}:\\
    $$
    \begin{aligned}
      & \text{Total children (N) = 4050} \\
      & \text{Unvaccinated (M) = 514} \\
      & \text{Vaccinated (N - M) = 3536} \\
      & \text{Sample size (n) = 65} \\
      \\
      & \text{Hypergeometric Distribution Formula:} \\
      & \boxed{P(X = k) = \frac{\binom{514}{k}
      \binom{3536}{65-k}}{\binom{4050}{65}}, \quad k = 0, 1, \dots, 65}
    \end{aligned}$$

  \item[27.] A display case contains thirty-five gems, of which ten
    are real diamonds and twenty-five are fake diamonds. A burglar
    removes four gems at random, one at a time and without
    replacement. What is the probability that the last gem she steals
    is the second real diamond in the set of four?

    \underline{Sol}:\\
    The burglar has to first draw a single diamond in the first set
    of three gems removed. Then, the burglar has to draw a second
    diamond on his fourth draw. We can calculate the probability of
    this to be \(\boxed{0.1289152}\)

  \item[28.] Consider an urn with r red balls and w white balls,
    where \(r + w = N\). Draw n balls in order without replacement.
    Show that the probability of k red balls is hypergeometric.

    \underline{Sol}:\\
    $$
    \begin{aligned}
      & \text{Total balls: } N = r + w \\
      & \text{Number of draws: } n \\
      & \text{Number of red balls desired: } k \\
      \\
      & \text{Total number of ways to choose } n \text{ balls from }
      N \text{ (ignoring order):} \\
      & \binom{N}{n} \\
      \\
      & \text{Number of ways to choose } k \text{ red balls from } r \text{:} \\
      & \binom{r}{k} \\
      \\
      & \text{Number of ways to choose the remaining } (n-k) \text{
      white balls from } w \text{:} \\
      & \binom{w}{n-k} \\
      \\
      & \text{By the Multiplication Principle, the number of
      successful outcomes is:} \\
      & \binom{r}{k} \binom{w}{n-k} \\
      \\
      & \text{The probability of } k \text{ red balls is the ratio of
      favorable to total outcomes:} \\
      & P(X = k) = \frac{\binom{r}{k} \binom{w}{n-k}}{\binom{N}{n}},
      \quad \max(0, n-w) \le k \le \min(n, r) \\
      \\
      & \text{This is the PMF of the Hypergeometric Distribution.}
    \end{aligned}$$

  \item[30.] Show that the ratio of two successive hypergeometric
    probability terms satisfies the following equation,
    \[
      \frac{\binom{r}{k + 1} \binom{w}{n - k - 1}}{\binom{N}{n}} \div
      \frac{\binom{r}{k} \binom{w}{n-k}}{\binom{N}{n}} = \frac{n -
      k}{k + 1} \cdot \frac{r - k}{w - n + k + 1}
    \]

    \underline{Sol}:\\
    $$
    \begin{aligned}
      & \text{Let } P(k) = \frac{\binom{r}{k}
      \binom{w}{n-k}}{\binom{N}{n}} \text{ and } P(k+1) =
      \frac{\binom{r}{k + 1} \binom{w}{n - k - 1}}{\binom{N}{n}}. \\
      \\
      & \frac{P(k+1)}{P(k)} = \frac{\binom{r}{k+1}
      \binom{w}{n-k-1}}{\binom{N}{n}} \cdot
      \frac{\binom{N}{n}}{\binom{r}{k} \binom{w}{n-k}} =
      \frac{\binom{r}{k+1}}{\binom{r}{k}} \cdot
      \frac{\binom{w}{n-k-1}}{\binom{w}{n-k}} \\
      \\
      & \text{Expanding the first term: } \\
      & \frac{\binom{r}{k+1}}{\binom{r}{k}} =
      \frac{\frac{r!}{(k+1)!(r-k-1)!}}{\frac{r!}{k!(r-k)!}} =
      \frac{k!(r-k)!}{(k+1)!(r-k-1)!} = \frac{k!}{(k+1)k!} \cdot
      \frac{(r-k)(r-k-1)!}{(r-k-1)!} = \frac{r-k}{k+1} \\
      \\
      & \text{Expanding the second term: } \\
      & \frac{\binom{w}{n-k-1}}{\binom{w}{n-k}} =
      \frac{\frac{w!}{(n-k-1)!(w-(n-k-1))!}}{\frac{w!}{(n-k)!(w-(n-k))!}}
      = \frac{(n-k)!(w-n+k)!}{(n-k-1)!(w-n+k+1)!} \\
      & \frac{\binom{w}{n-k-1}}{\binom{w}{n-k}} =
      \frac{(n-k)(n-k-1)!}{(n-k-1)!} \cdot
      \frac{(w-n+k)!}{(w-n+k+1)(w-n+k)!} = \frac{n-k}{w-n+k+1} \\
      \\
      & \text{Multiplying the simplified terms: } \\
      & \frac{P(k+1)}{P(k)} = \frac{r-k}{k+1} \cdot
      \frac{n-k}{w-n+k+1} = \frac{n - k}{k + 1} \cdot \frac{r - k}{w
      - n + k + 1}
    \end{aligned}$$

  \item[36.] Sixteen students: five freshmen, four sophomores, four
    juniors, and three seniors—have applied for membership in their
    school’s Communications Board, a group that oversees the
    college’s newspaper, literary magazine, and radio show. Eight
    positions are open. If the selection is done at random, what is
    the probability that each class gets two representatives? (Hint:
    Use the generalized hypergeometric model asked for in Question 3.2.35.)

    \underline{Sol}:\\
    Using question 3.2.35 as reference, we would get
    \[
      \frac{\binom{5}{2} \binom{4}{2} \binom{4}{2}
      \binom{3}{2}}{\binom{16}{8}} \approx \boxed{0.0839}
    \]
\end{enumerate}

\subsection*{3.3 Discrete Random Variables}
\begin{enumerate}
  \item[2.] An urn contains five balls numbered 1 to 5. Two balls are
    drawn simultaneously with replacement.
    \begin{enumerate}
      \item Let \(X\) be the larger of the two numbers drawn. Find \(p_X (k)\).

        \underline{Sol}:\\

        There are $5 \times 5 = 25$ total possible outcomes $(i, j)$.
        $X = \max(i, j)$.

        $$
        \begin{aligned}
          & \text{The event } X \le k \text{ occurs if both balls are
          } \le k. \\
          & P(X \le k) = \frac{k^2}{25} \\
          \\
          & p_X(k) = P(X \le k) - P(X \le k-1) \\
          & p_X(k) = \frac{k^2 - (k-1)^2}{25} = \frac{2k - 1}{25},
          \quad k = 1, 2, 3, 4, 5 \\
          \\
          & \text{Distribution Table for } X: \\
          &
          \begin{array}{c|ccccc}
            k & 1 & 2 & 3 & 4 & 5 \\
            \hline
            p_X(k) & \frac{1}{25} & \frac{3}{25} & \frac{5}{25} &
            \frac{7}{25} & \frac{9}{25} \\
          \end{array}
        \end{aligned}$$
      \item Let \(V\) be the sum of the two numbers drawn. Find \(p_V (k)\).

        \underline{Sol}:\\
        $V = i + j$, where $i, j \in \{1, 2, 3, 4, 5\}$. The possible
        sums range from $1+1=2$ to $5+5=10$.

        $$
        \begin{aligned}
          & \text{Outcomes for each sum } k: \\
          & k=2: \{(1,1)\} \implies 1 \\
          & k=3: \{(1,2), (2,1)\} \implies 2 \\
          & k=4: \{(1,3), (2,2), (3,1)\} \implies 3 \\
          & k=5: \{(1,4), (2,3), (3,2), (4,1)\} \implies 4 \\
          & k=6: \{(1,5), (2,4), (3,3), (4,2), (5,1)\} \implies 5 \\
          & k=7: \{(2,5), (3,4), (4,3), (5,2)\} \implies 4 \\
          & k=8: \{(3,5), (4,4), (5,3)\} \implies 3 \\
          & k=9: \{(4,5), (5,4)\} \implies 2 \\
          & k=10: \{(5,5)\} \implies 1 \\
          \\
          & \text{Distribution Table for } V: \\
          &
          \begin{array}{c|ccccccccc}
            k & 2 & 3 & 4 & 5 & 6 & 7 & 8 & 9 & 10 \\
            \hline
            p_V(k) & \frac{1}{25} & \frac{2}{25} & \frac{3}{25} &
            \frac{4}{25} & \frac{5}{25} & \frac{4}{25} & \frac{3}{25}
            & \frac{2}{25} & \frac{1}{25} \\
          \end{array}
        \end{aligned}$$
    \end{enumerate}

  \item[8.] Suppose a particle moves along the x-axis beginning at 0.
    It moves one integer step to the left with some probability. It
    is twice as likely to move to the right as to the left. What is
    the pdf of its position after four steps?

    \underline{Sol}:\\
    $$
    \begin{aligned}
      & \text{Let } p \text{ be the probability of moving left (} -1
      \text{) and } q \text{ be the probability of moving right (} +1
      \text{).} \\
      & q = 2p \\
      & p + q = 1 \implies p + 2p = 1 \implies p = \frac{1}{3}, \quad
      q = \frac{2}{3} \\
      \\
      & \text{Let } X \text{ be the position after } n=4 \text{
      steps. Let } k \text{ be the number of steps to the right.} \\
      & \text{The number of steps to the left is } 4 - k. \\
      & \text{Position } X = k - (4 - k) = 2k - 4. \\
      \\
      & \text{Possible values for } k \in \{0, 1, 2, 3, 4\}: \\
      & k=0 \implies X = -4 \\
      & k=1 \implies X = -2 \\
      & k=2 \implies X = 0 \\
      & k=3 \implies X = 2 \\
      & k=4 \implies X = 4 \\
      \\
      & \text{Probability mass function } p_X(x) = \binom{4}{k} q^k p^{4-k}: \\
      & p_X(-4) = \binom{4}{0} \left(\frac{2}{3}\right)^0
      \left(\frac{1}{3}\right)^4 = 1 \cdot 1 \cdot \frac{1}{81} =
      \frac{1}{81} \\
      & p_X(-2) = \binom{4}{1} \left(\frac{2}{3}\right)^1
      \left(\frac{1}{3}\right)^3 = 4 \cdot \frac{2}{3} \cdot
      \frac{1}{27} = \frac{8}{81} \\
      & p_X(0) = \binom{4}{2} \left(\frac{2}{3}\right)^2
      \left(\frac{1}{3}\right)^2 = 6 \cdot \frac{4}{9} \cdot
      \frac{1}{9} = \frac{24}{81} \\
      & p_X(2) = \binom{4}{3} \left(\frac{2}{3}\right)^3
      \left(\frac{1}{3}\right)^1 = 4 \cdot \frac{8}{27} \cdot
      \frac{1}{3} = \frac{32}{81} \\
      & p_X(4) = \binom{4}{4} \left(\frac{2}{3}\right)^4
      \left(\frac{1}{3}\right)^0 = 1 \cdot \frac{16}{81} \cdot 1 =
      \frac{16}{81} \\
      \\
      & \text{PDF for } X: \\
      &
      \begin{array}{c|ccccc}
        x & -4 & -2 & 0 & 2 & 4 \\
        \hline
        p_X(x) & \frac{1}{81} & \frac{8}{81} & \frac{24}{81} &
        \frac{32}{81} & \frac{16}{81} \\
      \end{array}
    \end{aligned}$$

  \item[14.] At the points \(x = 0, 1, \dots, 6 \), the cdf for the
    discrete random variable \(X\) has the value \(F_X = x(x +
    1)/42\). Find the pdf for \(X\).

    \underline{Sol}:\\
    $$
    \begin{aligned}
      & \text{The pdf } p_X(k) \text{ is related to the cdf } F_X(k)
      \text{ by:} \\
      & p_X(k) = F_X(k) - F_X(k-1) \\
      \\
      & \text{Given } F_X(x) = \frac{x(x+1)}{42} \text{ for } x = 0,
      1, \dots, 6: \\
      \\
      & p_X(k) = \frac{k(k+1)}{42} - \frac{(k-1)((k-1)+1)}{42} \\
      & p_X(k) = \frac{k^2 + k - (k-1)k}{42} \\
      & p_X(k) = \frac{k^2 + k - (k^2 - k)}{42} \\
      & p_X(k) = \frac{2k}{42} = \frac{k}{21}, \quad k = 1, 2, 3, 4, 5, 6 \\
      \\
      & \text{Distribution Table for } X: \\
      &
      \begin{array}{c|cccccc}
        k & 1 & 2 & 3 & 4 & 5 & 6 \\
        \hline
        p_X(k) & \frac{1}{21} & \frac{2}{21} & \frac{3}{21} &
        \frac{4}{21} & \frac{5}{21} & \frac{6}{21} \\
      \end{array}
    \end{aligned}$$
\end{enumerate}
